\chapter{Citations and Bibliographies}
\label{chap:citations}
\index{citation}\index{bibliography}\index{Bib\LaTeX{}}\index{Biber}\index{DOI}

A citation links a statement to a source; a reference gives the reader enough
information to identify that source. \LaTeX{} can format and renumber citations,
but it cannot decide whether a source supports a claim or whether its metadata
is true. We therefore begin with responsibility and data before software.

\begin{learningoutcomes}
  \item Distinguish citations, references, bibliographies, keys, and styles.
  \item Create verified records in a \file{.bib} database.
  \item Use \pkg{biblatex} with Biber to cite and print a bibliography.
  \item Produce textual, parenthetical, multiple, and page-specific citations.
  \item Recognise records for articles, books, chapters, datasets, software,
        theses, reports, and websites.
  \item Explain when a publisher's BibTeX or \pkg{natbib} workflow must be kept.
  \item Apply responsible practices for sources, permissions, and AI assistance.
\end{learningoutcomes}

\section{Concepts before commands}

A \term{citation} is the marker in the prose, such as an author and year or a
number. A \term{reference} is the fuller bibliographic description. A
\term{bibliography} is the organised set of references printed by the
document. A \term{bibliography database} is the separate source file, commonly
\file{references.bib}, that stores records.

Each record has a \term{citation key}, a private source identifier such as
\code{lamport1994}. The key is not printed. A \term{citation style} controls
the visible markers and bibliography. A DOI is a persistent identifier, not a
complete reference and not a guarantee that a source supports a claim.

\section{One verified record}

A Bib\LaTeX{} record starts with an entry type, key, and fields:

\begin{latexcode}{references.bib: a genuine book record}
@book{lamport1994,
  author    = {Lamport, Leslie},
  title     = {\LaTeX{}: A Document Preparation System},
  edition   = {2},
  publisher = {Addison-Wesley},
  date      = {1994},
  isbn      = {9780201529838}
}
\end{latexcode}

The record was checked against the publisher and the \LaTeX{} Project's book
list. The key could be different, but it must remain unique. Family-name-first
input helps Bib\LaTeX{} parse names. Do not guess a DOI because a book title looks
familiar.

\begin{goodpractice}
Open the source record or publisher page and verify author order, title,
container, year, volume, issue, pages, edition, and identifiers. Reference-
manager exports are starting data, not proof of accuracy.
\end{goodpractice}

\section{Connect the database with biblatex}

In the preamble, load \pkg{biblatex}, select Biber as the backend, and add the
database resource:

\begin{latexcode}{Complete example: Bib\LaTeX{} and Biber}
\documentclass{article}
\usepackage[backend=biber,style=authoryear]{biblatex}
\addbibresource{references.bib}

\begin{document}
\textcite{lamport1994} provides a recognised \LaTeX{} reference.
Structured source is also discussed in literate programming
\parencite{knuth1984}.

\printbibliography
\end{document}
\end{latexcode}

This complete project needs the two genuine records in
\file{references.bib}. \cmd{addbibresource} declares the database.
\cmd{textcite} places the author naturally in the sentence;
\cmd{parencite} produces a parenthetical citation. \cmd{printbibliography}
prints cited records in the selected style.

The reliable build command remains:

\begin{terminalcode}{Automated bibliography build}
latexmk -pdf main.tex
\end{terminalcode}

Internally, pdf\LaTeX{} writes citation requests, Biber reads the database and
creates formatted data, and pdf\LaTeX{} incorporates it on later passes. The
current Bib\LaTeX{} manual documents the data model and citation commands
\parencite{biblatexManual}.

\section{Useful citation forms}

\begin{latexcode}{Citation patterns}
\textcite{lamport1994} explains the document system.

The source-and-output distinction is longstanding
\parencite{lamport1994}.

Several source types can be cited together
\parencite{knuth1984,fisherIris}.

One passage can be identified precisely
\parencite[see][pp. 97--99]{knuth1984}.
\end{latexcode}

The first optional note appears before the citation and the second after it.
Page-specific citations help a reader verify a focused claim in a paginated
source. Do not add a page to a web page that has no stable pages.

\begin{pausepredict}
If a style changes from author--year to numeric, which part should remain
stable: the citation keys in source or the visible markers? This separation is
why keys should describe records rather than their temporary numbers.
\end{pausepredict}

\section{Records for scholarly source types}

The project database contains genuine examples of each major type:

\begin{itemize}
  \item journal article: \textcite{knuth1984};
  \item book: \textcite{lamport1994};
  \item chapter in a book: \textcite{chiroma2023};
  \item dataset: \textcite{fisherIris};
  \item software: \textcite{rproject};
  \item doctoral thesis: \textcite{heffernan2020};
  \item technical standard/report: \textcite{iso5725}; and
  \item website/documentation: \textcite{latexprojectDocumentation}.
\end{itemize}

Entry types select suitable fields and formatting. For software used in an
analysis, obtain the citation from that installed version when possible; the R
documentation recommends its \code{citation()} function. For datasets, cite
the version and persistent identifier supplied by the repository. For a web
resource that changes, record an access date.

\section{Common bibliography failures}

\begin{commonerror}
\textbf{Symptom:} citation keys appear as bold text or question marks and the
bibliography is absent.

\textbf{Checks:} confirm the key exists and matches case; confirm the database
path; read the Biber log; make sure \cmd{printbibliography} is present; and run
the correct build sequence. An earlier \LaTeX{} error can prevent a valid control
file from reaching Biber.
\end{commonerror}

A warning about an empty bibliography can mean that nothing was cited. Use
\cmd{nocite}\required{*} only when the document intentionally needs every
database record, not as a way to hide missing citations.

Special characters in names and titles need correct Bib\LaTeX{} encoding. Protect
capital letters whose case must remain, such as an acronym, with braces after
checking the style behaviour. Do not wrap every title in extra braces; that can
defeat style-controlled case conversion.

\section{Publisher workflows: BibTeX and natbib}

Many publishers supply older but valid workflows built around BibTeX,
\pkg{natbib}, and \file{.bst} bibliography-style files. BibTeX reads
\file{.bib} data and uses the \file{.bst} file to format a \file{.bbl};
\pkg{natbib} supplies citation commands for many author--year and numerical
styles.

\begin{submissioncheck}
If a journal template uses BibTeX or \pkg{natbib}, follow its current
instructions. Do not load \pkg{biblatex} alongside it or replace its
bibliography system merely because this chapter uses Bib\LaTeX{} for teaching.
Submit the \file{.bbl} when the journal requests it.
\end{submissioncheck}

The educational default and a publisher requirement can legitimately differ.
The important skill is recognising the existing workflow and building it in
the documented sequence.

\section{Scholarly responsibility}

A citation must support the nearby statement. Citing a paper after reading only
its title or an AI summary is not sufficient. Read the relevant source, prefer
primary evidence for a primary claim, represent uncertainty honestly, and do
not cite a review as though it performed every study it summarises.

Quotation requires exactness and sometimes permission. Figures, questionnaires,
and substantial table content can require permissions even when cited. Data
and software are research contributions and should be credited using the
records their creators request.

AI-assisted tools may help search, reformat, or inspect a draft record, but they
can invent authors, titles, dates, and DOIs. The author remains responsible for
checking each record and claim, complying with disclosure rules, protecting
confidential information, and following publisher or institutional policy.

\begin{scienceinpractice}
The Iris dataset record identifies a repository DOI and an explicit dataset
type \parencite{fisherIris}. That citation acknowledges the data source; it
does not automatically justify using the dataset for a new scientific claim.
The methods must still explain selection, processing, and limitations.
\end{scienceinpractice}

\section{Code lab: build a bibliography from two files}

\begin{codelab}
Create \path{main.tex} and \path{references.bib} in the same folder. Run
\pkg{latexmk} as shown. Then change the citation key in only the source file,
observe the warning, and restore the matching key.
\end{codelab}

\begin{latexcode}{main.tex}
\documentclass{article}
\usepackage[backend=biber,style=authoryear]{biblatex}
\addbibresource{references.bib}

\begin{document}
The system is maintained by the \textcite{latexproject}.
Its documentation is available online \parencite{latexproject}.
\printbibliography
\end{document}
\end{latexcode}

\begin{latexcode}{references.bib}
@online{latexproject,
  author  = {{The \LaTeX{} Project}},
  title   = {The \LaTeX{} Project},
  url     = {https://www.latex-project.org/},
  urldate = {2026-07-31}
}
\end{latexcode}

\begin{terminalcode}{Build the document and its bibliography}
latexmk -pdf main.tex
\end{terminalcode}

\chapterproject

\begin{miniproject}
Create \file{references.bib} for the running article. Add the verified ISO
measurement-method standard and one verified \LaTeX{}/software record used by the
workflow. Load \pkg{biblatex} with Biber, cite sources only where they support
the prose, and print the bibliography. Run \pkg{latexmk}; correct every
undefined citation before continuing.
\end{miniproject}

\chaptersummary

Citations connect claims to references stored as verified database records.
Keys are stable source identifiers; styles control visible form. Bib\LaTeX{} and
Biber support structured entry types and varied citation commands.
\pkg{latexmk} automates the required passes. Publisher BibTeX/\pkg{natbib}
workflows must be preserved when required. Accuracy, relevance, permission,
data/software credit, and responsible AI use remain the author's obligations.

\chaptercheck
\begin{chapterchecklist}
  \item Every record was checked against an authoritative source.
  \item Citation keys are unique, stable, and meaningful.
  \item Each citation supports the claim beside it.
  \item Data and software used in the work receive appropriate credit.
  \item The selected workflow matches the journal or institutional instruction.
  \item No undefined citation or empty bibliography warning remains.
\end{chapterchecklist}

\chapterexercisestitle
\begin{chapterexercises}
  \item Distinguish a citation key, citation, reference, and bibliography.
  \item Create a verified \code{@article} record from a publisher page.
  \item Produce textual and parenthetical citations to one record.
  \item Cite two records together and add a page-specific postnote.
  \item Diagnose a missing bibliography caused by an absent print command.
  \item Explain why a DOI alone is not evidence that a source supports a claim.
  \item State when to keep a publisher's BibTeX/\pkg{natbib} workflow.
  \item Challenge: audit five records for metadata accuracy and document the
        authoritative page used for each check.
\end{chapterexercises}
